\scanpage{328}

so the condition reads $\unsure{v|\nu}=-1$, i.e. this is condition $3^\circ$, but then the expression (23) for $\rho^3$ shows us that $\rho^3=-1$.

\uline{\unsure{Remarks|Remark}} The preceding conditions imply the formulas
\[
\begin{aligned}
(24)\qquad &\varepsilon_0=\sigma\rho\\
(25)\qquad &\varepsilon_1=\rho\sigma\\
(26)\qquad &\rho\varepsilon_1=\varepsilon_0^{-1}\\
(27)\qquad &\rho^2\sigma=\varepsilon_0^{-1}
\end{aligned}
\]
each of which \inserted{except the last (27)} implies them in turn
--- for example we have $\varepsilon_0=-\sigma^{-1}\rho$ \inserted{(7)}, so (24) is equivalent to \inserted{$\sigma=-\sigma^{-1}$ i.e.} $\sigma^2=-1$, and (25) is equivalent to (24) by applying $\operatorname{int}(\rho)$ to both \unsure{sides}.

We have moreover
\[
\rho\varepsilon_1=
\begin{pmatrix}0&1\\\unsure{v|\nu}&\unsure{w|\omega}\end{pmatrix}
\begin{pmatrix}1&0\\1&1\end{pmatrix}
=\begin{pmatrix}1&1\\\unsure{v|\nu}+\unsure{w|\omega}&\unsure{\unsure{w|\omega}^{\prime}|\unsure{w|\omega}}\end{pmatrix},
\qquad
\varepsilon_0^{-1}=\begin{pmatrix}1&1\\0&1\end{pmatrix}
\]
which shows that $\rho\varepsilon_1=\varepsilon_0^{-1}$ is equivalent to:
\[
\unsure{w|\omega}=1,\quad \unsure{v|\nu}+\unsure{w|\omega}=0\qquad\text{i.e.}\quad \unsure{w|\omega}=1,\unsure{v|\nu}=-1\quad\text{i.e.}\quad
\rho=\begin{pmatrix}0&1\\-1&1\end{pmatrix},
\]
OK. Composing (25) \inserted{(in the form $\rho\varepsilon_1=\rho^2\sigma$)} and (26) we find (27) as a consequence of the conditions $2^\circ)+3^\circ)$; \deleted{\unsure{the}}\inserted{Let us try to prove} converse implication \deleted{\unsure{follows}}\inserted{by} \deleted{from} a brute-force calculation
\[\deleted{\rho^2\sigma=}\rho^2\sigma=
\begin{pmatrix}
-\unsure{v|\nu}\unsure{w|\omega}&\unsure{w|\omega}-(\unsure{v|\nu}+\unsure{w|\omega}^2)\\
-\unsure{v|\nu}(\unsure{v|\nu}+\unsure{w|\omega}^2)&\unsure{-\unsure{w|\omega}(\unsure{w|\omega}^2+2\unsure{v|\nu})+\unsure{v|\nu}-\unsure{w|\omega}^2|-\unsure{w|\omega}(\unsure{w|\omega}^2+2\unsure{v|\nu})+\unsure{v|\nu}+\unsure{w|\omega}^2}
\end{pmatrix},
\qquad
\varepsilon_0^{-1}=\begin{pmatrix}1&\unsure{-1/\unsure{v|\nu}}\\0&1\end{pmatrix}
\]
\deleted{since $\unsure{v|\nu}$ invertible} we can have the equality $\rho^2\sigma=\unsure{-\varepsilon_0^{-1}|\varepsilon_0^{-1}}$ only if $-\unsure{v|\nu}(\unsure{v|\nu}+\unsure{w|\omega}^2)=0$, \unsure{this implies|which implies} ($\unsure{v|\nu}$ invertible!) $\unsure{v|\nu}+\unsure{w|\omega}^2=0$, \deleted{\unsure{whence}}\inserted{i.e. $\unsure{v|\nu}=-\unsure{w|\omega}^2$},

\mspage{327}

\scanpage{329}
\mspage{465}

\[
\rho^2\sigma=\begin{pmatrix}\unsure{w|\omega}^3&\unsure{w|\omega}\\0&\deleted{\unsure{-w(1+2w)}}\inserted{\unsure{w|\omega}^3}\end{pmatrix},
\]
This shows that
\deleted{$\rho^2\sigma=\varepsilon_0^{-1}$ \unsure{iff} $\unsure{w|\omega}=1$ OK.}
\margin{\deleted{\unsure{We have in addition to} \illegible{4}}
\unsure{$\unsure{v|\nu}=-\unsure{w|\omega}^2$} \deleted{\illegible{2}}
\unsure{i.e. $\unsure{v|\nu}=-1,\unsure{w|\omega}=1$}}
\inserted{
\[
(27\mathrm{bis})\qquad
\rho^2\sigma=\varepsilon_0^{-1}\Longleftrightarrow
\left\{\unsure{v|\nu}+\unsure{w|\omega}^2=0,\ \unsure{w|\omega}^3=1\right.
\quad
\left(\Rightarrow \unsure{v|\nu}=-\unsure{w|\omega}^2=-1/\unsure{w|\omega}\ \text{i.e.}\ \unsure{w|\omega}=-1/\unsure{v|\nu},\ \mathrm{OK}\right)
\]
}

But (27), on multiplying by $\rho$, is also equivalent to
\[\deleted{\unsure{(27)}\,\illegible{1}}\rho^3\sigma=\deleted{\unsure{u}}\rho\varepsilon_0^{-1}
\]
now we know that $\varepsilon_0=\unsure{-\sigma^{-1}\rho}$ \inserted{(7)} i.e. $\varepsilon_0^{-1}=-\deleted{\unsure{\sigma}}\rho^{-1}\unsure{\sigma|\sigma_0}$
i.e. $\rho\varepsilon_0^{-1}=-\sigma$, so the preceding relation also reads $\rho^3\sigma=-\sigma$ or again
\[
(28)\qquad \boxed{\rho^3=-1}.
\]
We can see \deleted{\unsure{this}} the equivalence \deleted{directly}\inserted{of (28) with the conditions (27bis)} by computing $\rho^3$ in terms of $\unsure{v|\nu},\unsure{w|\omega}$, we find
\[
(29)\qquad
\rho^3=\begin{pmatrix}\unsure{v|\nu}\unsure{w|\omega}&\unsure{v|\nu}+\unsure{w|\omega}^2\\\unsure{v|\nu}(\unsure{v|\nu}+\unsure{w|\omega}^2)&\unsure{w|\omega}(2\unsure{v|\nu}+\unsure{w|\omega}^2)\end{pmatrix}
\]
And we can have $\widetilde\rho^3=1$ i.e. $\rho^3$ a scalar only if \deleted{$\unsure{v|\nu}+\unsure{w|\omega}^2=$}
\[
\unsure{v|\nu}+\unsure{w|\omega}^2=0\qquad\text{i.e.}\quad \unsure{v|\nu}=-\unsure{w|\omega}^2
\]
which implies
\[
\rho^3=\begin{pmatrix}-\unsure{w|\omega}^3&0\\0&-\unsure{w|\omega}^3\end{pmatrix}
\]
i.e. we find

\uline{Lemma} We have $\widetilde\rho^3=1$ iff $\unsure{v|\nu}+\unsure{w|\omega}^2=0$ i.e. $\unsure{v|\nu}=-\unsure{w|\omega}^2$ i.e. we have
\[
(30)\qquad
\rho=\begin{pmatrix}0&1\\-\unsure{w|\omega}^2&\unsure{w|\omega}\end{pmatrix},
\]
and then we have
\[
(31)\qquad \rho^3=\unsure{-}\unsure{w|\omega}^3\cdot1\ ,
\]
so we have
\[
(32)\qquad
\rho^3=-1\Longleftrightarrow
\deleted{\unsure{\unsure{v|\nu}=-\unsure{w|\omega}^2}}\quad
\begin{gathered}
\unsure{w|\omega}^3=1\ \text{i.e. }\unsure{w|\omega}\in\Gamma\mu_3\\
\inserted{\uline{and}\quad \unsure{v|\nu}=-\unsure{w|\omega}^2\ (= -1/\unsure{w|\omega})}
\end{gathered}
\]

\mspage{328}

\scanpage{330}

\uline{Corollary 1} The equivalent conditions (22) are also equivalent to
\[
(33)\qquad \rho^3=-1\quad\text{and}\quad\det\rho=1
\]
or again to
\[
(34)\qquad \rho^3=-1\quad\text{and}\quad\deleted{\unsure{\det}}\inserted{\operatorname{Tr}}\rho=1
\]
in other words, considering the three relations
\[
(35)\qquad \det\rho=1,\quad\operatorname{Tr}\rho=1,\quad\rho^3=-1
\]
\inserted{the conjunction of} two of these \unsure{imply|implies} the third, and is equivalent to the equivalent conditions of (22).

\margin{false}
\deleted{
Let us note another way of finding this corollary, by noting that
\[
(\unsure{36|30})\qquad
\rho^3+1=(\rho+1)(\rho^2-\rho+1)\qquad
\text{(algebraic \uline{identity})}
\]
now $\rho$ not being a scalar on any fiber,
\uline{$\rho+1$ is invertible}, so $\rho^3=-1$ i.e. $\rho^3+1=0$ is equivalent to the relation $\rho^2-\rho+1=0$:
\[
(36)\qquad \rho^2-\rho+1=0\ ;\quad(\Longleftrightarrow\rho^3=-1),
\]
which we compare with Hamilton--Cayley
\[
(37)\qquad
\rho^2-\underbrace{(\operatorname{Tr}\rho)}_{\unsure{w|\omega}}\rho+\underbrace{\det\rho}_{-\unsure{v|\nu}}=0,
\]
so \deleted{$\unsure{w|\omega}=1,\unsure{v|\nu}=-1$} $\operatorname{Tr}\rho=\det\rho=1$ implies (36), \deleted{so $\rho^3=-1$}
on the other hand if we have (36), we conclude via 37
\[
(\operatorname{Tr}(\rho)-1)\rho\ \unsure{=|\ne}\ \deleted{\unsure{\det}}\,(\det\rho-1)1
\]
}

\mspage{329}

\scanpage{331}
\mspage{466}

Another way of finding certain equivalences is via the Hamilton--Cayley equations for $\sigma$ and $\rho$
\[
(36)\qquad
\deleted{\unsure{\rho}}\inserted{\sigma}^2
-\underbrace{(\operatorname{Tr}\sigma)}_{\unsure{w|\omega}'=1-\unsure{w|\omega}}\sigma
+\underbrace{\det\sigma}_{\unsure{v|\nu}'=-\unsure{v|\nu}}=0
\]
\[
(37)\qquad
\rho^2-\underbrace{(\operatorname{Tr}\rho)}_\unsure{w|\omega}\rho
+\underbrace{\det\rho}_{-\unsure{v|\nu}}=0.
\]
Equation (36) gives again that the conjunction of $\operatorname{Tr}\sigma=0$, $\det\sigma=1$ implies $\sigma^2+1=0$ i.e. $\sigma^2=-1$, and conversely, taking into account that $1$ and $\sigma$ are \inserted{sections} \unsure{linear\textsuperscript{ly}} independent \unsure{in|of} $\unsure{\mathrm{M}_2}(\unsure{\underline{O}_S|\mathcal O_S})$.
\unsure{Likewise|=} (37) shows
\[
(38)\qquad
(\operatorname{Tr}\rho=1\ \text{and}\ \det\rho=1)
\Longleftrightarrow(\rho^2-\rho+1=0).
\]
and we \deleted{relation} have the implication
\[
(39)\qquad
\rho^2-\rho+1=0\quad\Longrightarrow\quad
\rho^3+1=0\ \text{i.e.}\ \rho^3=-1
\]
by virtue of the \uline{algebraic} identity
\[
\rho^3+1=(\rho+1)(\rho^2-\rho+1).
\]
So

\uline{Corollary 2}. The equivalent conditions of (22) are also equivalent to the relation
\[
(40)\qquad \rho^2+\rho+1=0.
\]

\mspage{330}

\scanpage{332}

\deleted{
$5^\circ)$ Modular pinning theorem.

\uline{Theorem} Let $G,SG,PG$ be forms \inserted{(over scheme $S$)} of $\mathrm{Gl}(2),\mathrm{Sl}(2),\mathrm{GP}(1)$ \inserted{\deleted{\unsure{over scheme}}} \unsure{(}associated\unsure{)}, and $\widetilde\sigma,\widetilde\rho$ two sections of $PG$.
}

$5^\circ)$ \uline{A return to} $\uline{1^\circ)}$, concerning the situation
\[
(41)\quad
\left\{\begin{array}{l}
\widetilde\rho,\widetilde\sigma\ \text{with }1^\circ)\quad
\widetilde\varepsilon_0=\widetilde\sigma^{-1}\widetilde\rho\quad
\uline{\text{regular}}\text{ unipotent}\ \deleted{\text{\unsure{i.e.}}}\quad(\text{i.e.}\\
\widetilde\varepsilon_1=\widetilde\rho\widetilde\sigma^{-1}
=\widetilde\rho(\widetilde\varepsilon_0)
=\widetilde\sigma(\widetilde\varepsilon_0)\quad
\uline{\text{regular}}\text{ unipotent})\ \text{and}\\
2^\circ)\quad\widetilde\varepsilon_0\ \text{and}\ \widetilde\varepsilon_1\ \text{inequivalent in each fiber.}
\end{array}\right.
\]

Let $\varepsilon_0,\varepsilon_1$ be the \uline{regular} unipotent liftings in $G$ --- identified with sections of $\unsure{M|\mathcal M}$
\[
(42)\qquad\varepsilon_0=1+N_0,\quad\varepsilon_1=1+N_1,
\]
\margin{$N_0$ and $N_1$ strictly \unsure{nilpoten\textsuperscript{t}} i.e. their determinant and their trace \deleted{\illegible{1}} zero}
with
\[
(43)\qquad
\widetilde\sigma(\varepsilon_0)=\widetilde\rho(\varepsilon_0)=\varepsilon_1,
\quad\inserted{\text{i.e.,}}\quad
\widetilde\sigma(N_0)=\widetilde\rho(N_0)=N_1.
\]
We can then lift $\widetilde\sigma$ and $\widetilde\rho$ canonically to sections $\sigma,\rho$ of $G$ (or of $\unsure{M|\mathcal M}$) satisfying the following conditions
\[
(44)\quad\left\{
\begin{array}{ll}
1^\circ)\ -\sigma^{-1}\rho=\varepsilon_0
&[\Longleftrightarrow-\rho\sigma^{-1}=\varepsilon_1]\\
\deleted{\unsure{2^\circ)(\rho+\sigma)}}
&\inserted{\left[\begin{aligned}
&\Longleftrightarrow\rho+\sigma=-N_1\sigma\ (= -\sigma N_0)\\
&\Longleftrightarrow\rho+\sigma=N_1\rho\ (=\rho N_0)
\end{aligned}\right]}\\
&\deleted{\unsure{\Longleftrightarrow\rho+\sigma=}}\\
2^\circ)\ (\rho+\sigma)^2=\rho+\sigma
&\inserted{\begin{gathered}
\Longrightarrow(\rho+\sigma)^2=\unsure{u}(\rho+\sigma)\quad\text{with}\\
u\in\Gamma G_m
\end{gathered}}
\end{array}\right.
\]
Moreover, there exists a \deleted{\unsure{isomorphism}} unique triple $(\varphi,\rho_0,\sigma_0)$
\[
(45)\quad\left\{\begin{array}{l}
G_0=\mathrm{Gl}(2)_S\xrightarrow{\varphi}G\quad(\text{pinning})\\
\rho_0,\sigma_0\in\Gamma(G_0),\quad
\rho_0=\begin{pmatrix}0&1\\\unsure{v|\nu}&\unsure{w|\omega}\end{pmatrix},\quad
\sigma_0=\begin{pmatrix}0&-1\\\unsure{v|\nu}'&\unsure{w|\omega}'\end{pmatrix}\\
(\unsure{v|\nu}+\unsure{v|\nu}'=0,\unsure{w|\omega}+\unsure{w|\omega}'=1),\quad \unsure{v|\nu}\in\Gamma G_m\\
\inserted{\text{i.e. }\unsure{v|\nu}'=-\unsure{v|\nu},\unsure{w|\omega}'=1-\unsure{w|\omega}}
\end{array}\right.
\]
such that we have
\[
(\unsure{46|48})\qquad\varphi(\rho_0)=\rho,\quad\varphi(\sigma_0)=\sigma,
\]
which implies
\[
(47)\qquad\varphi((\varepsilon_0)_0)=\varepsilon_0,\quad
\varphi((\varepsilon_1)_0)=\varepsilon_1
\]

\mspage{331}